\chapter{Parametri e taratura}
\label{ch:parametri}

\section{Quadro completo dei parametri}

Valori predefiniti della versione in uso, \code{highs\_survey\_prior}.

\begin{center}
\begin{tabularx}{\textwidth}{@{}l r X@{}}
\toprule
\textbf{Parametro} & \textbf{Default} & \textbf{Effetto} \\
\midrule
\multicolumn{3}{@{}l}{\itshape\color{CorpGray} Discretizzazione} \\
\code{granularity\_min} & 15 & ampiezza dello slot in minuti; principale leva sul costo computazionale \\
\code{resample\_method} & \code{mean} & funzione di aggregazione nel ricampionamento \\
\code{time\_limit} & 60 & budget del solver per ciascun giorno, in secondi \\
\addlinespace
\multicolumn{3}{@{}l}{\itshape\color{CorpGray} Modello di potenza} \\
\code{power\_level\_variation} ($\delta$) & 0.15 & scarto dei livelli basso e alto dal valore tipico \\
\addlinespace
\multicolumn{3}{@{}l}{\itshape\color{CorpGray} Penalità temporali} \\
\code{time\_window\_penalty} ($\lambda_w$) & 1.0 & peso della penalità di fascia oraria \\
\code{window\_penalty\_ramp\_min} ($\rho$) & 60 & minuti di distanza che aggiungono 1 al fattore \\
\code{window\_penalty\_max\_factor} ($F_{\max}$) & 6.0 & saturazione del termine di distanza \\
\code{window\_vagueness\_weight} ($w$) & 0.2 & peso del pavimento di genericità \\
\code{season\_penalty} ($\lambda_s$) & 8.0 & costo per slot fuori dai mesi dichiarati \\
\addlinespace
\multicolumn{3}{@{}l}{\itshape\color{CorpGray} Penalità di utilizzo} \\
\code{activation\_penalty} ($\lambda_a$) & 1.0 & costo per ogni accensione \\
\code{duration\_penalty\_block} ($\lambda_d$) & 0.5 & costo per slot di scostamento dalla durata \\
\code{duration\_penalty\_daily} ($\lambda_D$) & 0.0 & budget sul monte ore giornaliero (disattivato) \\
\code{over\_activation\_factor} & 1.5 & moltiplicatore del punto di pareggio della quota \\
\bottomrule
\end{tabularx}
\end{center}

\section{Punti di pareggio}

Riepilogo dell'interpretazione dei coefficienti, discussa in dettaglio nel
capitolo~\ref{ch:obiettivo}.

\begin{center}
\begin{tabularx}{\textwidth}{@{}l l X@{}}
\toprule
\textbf{Coefficiente} & \textbf{Pareggio} & \textbf{Significato del pareggio} \\
\midrule
$\lambda_w$ & 1.0 & un'attivazione fuori fascia costa quanto non spiegare $p_i$ per uno slot \\
$\lambda_a$ & 1.0 & un'accensione in più costa quanto uno slot acceso a segnale nullo \\
$\lambda_d$ & 1.0 & uno slot di sforamento costa quanto uno slot inspiegato \\
$\lambda_D$ & 1.0 & uno slot di scarto dal monte ore costa quanto uno slot inspiegato \\
$\lambda_{q,i}$ & $s_i$ & \textbf{eccezione}: il pareggio è la durata tipica del ciclo in slot \\
$\lambda_s$ & 1.0 & impostato a 8 di proposito, ben oltre il pareggio \\
\bottomrule
\end{tabularx}
\end{center}

\section{Osservazioni sperimentali}

\subsection{La penalità di fascia oraria è troppo forte al valore predefinito}

Con $\lambda_w = 1$ la penalità di fascia eguaglia esattamente il costo di lasciare il consumo
inspiegato. In queste condizioni il modello, indifferente fra le due opzioni, non attiva mai
l'apparecchio fuori fascia: il vincolo formalmente flessibile si comporta come rigido.

L'effetto è stato osservato direttamente in una prova controllata: con $\lambda_w = 1$ un
apparecchio non veniva attivato in alcun punto della giornata. Riducendo il coefficiente a $0.3$ il
comportamento diventa quello atteso --- attivazione dentro e appena fuori dalla fascia, rifiuto a
grande distanza.

\begin{notebox}{Raccomandazione}
Per i modelli con obiettivo $L_1$ si consiglia $\lambda_w$ nell'intervallo $0.2$--$0.4$. Il valore
predefinito di $1.0$ è ereditato dai modelli quadratici, dove la scala delle penalità è diversa, e
in questo contesto risulta eccessivo.
\end{notebox}
